\section{均衡不等于最优}
\label{sec:17-Constrained-Guarantees/06-Games-and-Equilibria/03}

一个工单系统上线了``跨部门直通''。原来的两条流转路径各有一段排队环节和一段固定审批，新加的这条连接几乎不占时间，本意是让积压的一侧可以借道另一侧。三个月后平均处理时长从 \(65\) 小时涨到 \(80\) 小时。没有人多做一步操作，没有一个环节变慢，加进去的那条通道本身耗时接近零。

上一节的结论已经排除了最容易想到的解释：这不需要有人做错决定。每个提单人都在按当时最快的路径走，每一步都是理性的。可以问的是，能不能把``所有人各自最优''与``整体最优''之间的差距算出来一个数。

\subsection{给这个差距一个比值}

先把上面的情形写成一个能算清的模型。这是 Braess 在 1968 年给出的例子，用道路网络叙述最直观。

起点 \(S\) 到终点 \(T\)，中间有两个节点 \(A\) 与 \(B\)。\(S\to A\) 这段是拥堵路段，通行时间等于该段车流量除以 \(100\)；\(A\to T\) 是一条不受流量影响的老路，固定 \(45\) 分钟。另一侧对称：\(S\to B\) 固定 \(45\) 分钟，\(B\to T\) 的通行时间是流量除以 \(100\)。早高峰有 \(4000\) 辆车要从 \(S\) 到 \(T\)。

均衡在两条路各走 \(2000\) 辆时达成：每辆车的通行时间是 \(2000/100+45=65\) 分钟，两条路一样长，谁换路都不会更快。

现在修一条 \(A\to B\) 的连接，通行时间为零。多出的路线是 \(S\to A\to B\to T\)，它的通行时间等于两段拥堵路段之和。对一辆车来说这条新路很有吸引力：只要 \(S\to A\) 上的车少于 \(4500\) 辆，走 \(S\to A\) 就比走 \(S\to B\) 那 \(45\) 分钟快；到了 \(A\) 之后同样的算式让它继续走 \(B\to T\) 而不是 \(A\to T\)。

于是所有 \(4000\) 辆车都挤上新路线，通行时间变成 \(4000/100+0+4000/100=80\) 分钟。

这是不是均衡，要逐条验。一辆车若改走 \(S\to A\to T\)：\(S\to A\) 上仍有 \(4000\) 辆车，\(40\) 分钟，加上 \(A\to T\) 的 \(45\) 分钟，共 \(85\) 分钟，比 \(80\) 慢。改走 \(S\to B\to T\)：\(45+40=85\) 分钟，同样慢。没有人能靠自己改路变快，\(80\) 分钟就是新的均衡。

修路之前 \(65\) 分钟，修了一条零耗时的路之后 \(80\) 分钟。所有人都变慢了，而且没有一个人做错了选择。

\begin{center}
\begin{tikzpicture}[x=1cm,y=1cm,font=\small,>=stealth]
  \foreach \dx in {0,5.9}{
    \begin{scope}[shift={(\dx,0)}]
      \draw[->,black!65] (0.22,1.0) -- (1.38,1.78);
      \draw[->,black!65] (1.82,1.78) -- (2.98,1.0);
      \draw[->,black!65] (0.22,0.8) -- (1.38,0.02);
      \draw[->,black!65] (1.82,0.02) -- (2.98,0.8);
      \draw[black!75,fill=white] (0,0.9) circle (0.22);
      \draw[black!75,fill=white] (1.6,1.9) circle (0.22);
      \draw[black!75,fill=white] (1.6,-0.1) circle (0.22);
      \draw[black!75,fill=white] (3.2,0.9) circle (0.22);
      \node[font=\scriptsize] at (0,0.9) {\(S\)};
      \node[font=\scriptsize] at (1.6,1.9) {\(A\)};
      \node[font=\scriptsize] at (1.6,-0.1) {\(B\)};
      \node[font=\scriptsize] at (3.2,0.9) {\(T\)};
      \node[font=\scriptsize,black!70,anchor=south east] at (0.85,1.44) {\(x/100\)};
      \node[font=\scriptsize,black!70,anchor=south west] at (2.35,1.44) {\(45\)};
      \node[font=\scriptsize,black!70,anchor=north east] at (0.85,0.36) {\(45\)};
      \node[font=\scriptsize,black!70,anchor=north west] at (2.35,0.36) {\(x/100\)};
    \end{scope}
  }
  \draw[->,blue!65!black,thick] (7.5,1.68) -- (7.5,0.12);
  \node[font=\scriptsize,blue!65!black,anchor=west] at (7.62,0.9) {\(0\)};
  \node[align=left,font=\footnotesize,anchor=north] at (1.6,-0.65)
    {加路前：两条路各 \(2000\) 辆\\ 每辆 \(20+45=65\) 分钟};
  \node[align=left,font=\footnotesize,anchor=north] at (7.5,-0.65)
    {加路后：\(4000\) 辆全走 \(S\!\to\!A\!\to\!B\!\to\!T\)\\ 每辆 \(40+0+40=80\) 分钟};
\end{tikzpicture}

\emph{右图只比左图多了一条通行时间为零的边，均衡通行时间却从 \(65\) 涨到 \(80\)。往系统里增加选项不是无害的。}
\end{center}

图里那条蓝色的边是唯一的改动。它不消耗任何时间，也不减少任何人的选项——每个人加路之后仍然可以走原来的路。可正是``可以走''让所有人都走了，而所有人都走的结果是所有人都变慢。这与直觉冲突的地方在于：单人优化里增加可行域绝不会让最优值变差，而这里的``最优''是均衡而不是最优化，增加可行域会改变均衡的位置。

把差距写成一个数。定义\emph{无政府代价}为最差 Nash 均衡的社会成本与社会最优成本之比：

\[
\mathrm{PoA}=\frac{\max_{s\in\mathrm{NE}}C(s)}{\min_{s}C(s)}.
\]

分子是让每个人自己选路的结果，分母是一个可以强制指派路线的调度者能达到的最好结果。加路之后的网络里，社会最优仍然是让 \(2000\) 辆走上路、\(2000\) 辆走下路、没人用新边，总时长对应每人 \(65\) 分钟，而均衡是 \(80\)，于是 \(\mathrm{PoA}=80/65\approx1.23\)。

这个比值有上界，而且很紧。对所有延迟函数为仿射形式 \(a x+b\)（\(a,b\ge0\)）的路网，Roughgarden 与 Tardos 证明 \(\mathrm{PoA}\le4/3\)，与网络拓扑无关。达到 \(4/3\) 的极端例子只有两条边：从 \(S\) 到 \(T\) 有一条延迟为 \(x\) 的拥堵路和一条固定 \(1\) 的路，总流量为 \(1\)。均衡时所有人走拥堵路（因为 \(x\le1\) 恒成立），总成本 \(1\)；最优是一半一半，总成本 \(\tfrac12\cdot\tfrac12+\tfrac12\cdot1=0.75\)；比值恰好 \(4/3\)。Braess 例子里的 \(1.23\) 已经很接近这个上界了。

上界的存在是个好消息：仿射延迟下自私选路最多浪费三分之一。同时要看清它有多脆。把延迟函数换成多项式，\(\mathrm{PoA}\) 随次数增长而增长；换成更一般的函数，它可以无界。所以``自私选路只差一点''这句话是有条件的定理，不是普遍规律。

\subsection{改规则，不是劝人改行为}

均衡由收益结构决定。既然如此，想让均衡换个位置，能动的就不是参与者的觉悟，而是收益结构本身。这是机制设计的全部出发点。

拿上面那个两条边的例子看具体怎么动。均衡之所以偏离最优，是因为多走一辆车会让\emph{其他所有人}都慢一点，而这份影响没有记在这辆车自己的账上。把它记上去就行：在拥堵路上按流量 \(f\) 收一笔等于边际外部影响的费用。拥堵路的总时长是 \(f\cdot f\)，对 \(f\) 求导得 \(2f\)，其中 \(f\) 是自己承受的，另外那个 \(f\) 是施加给别人的，所以收费定为 \(f\)。司机感知到的成本变成 \(f+f=2f\)，与固定路的 \(1\) 相等时 \(f=1/2\)——正好是社会最优的分流比例。

这个数不是凑出来的。它是社会成本对流量的边际值，也就是拥堵约束的影子价格，与\hyperref[sec:09-Convex-Optimization/05-Duality/04]{``对偶变量是约束的影子价格''}里那个量是同一个东西。伦敦、新加坡、斯德哥尔摩的拥堵收费方案背后都是这条推导。

同样的动作在别处换了名字。排放配额把污染的外部影响折进企业的成本函数；广告位拍卖用出价规则决定谁拿到位置以及付多少；平台对刷单的处罚把作弊那一格的收益压到低于诚实那一格。这些都不是在说服参与者，而是在重写矩阵。

拍卖里有一个说得清的例子。第二价格拍卖中，出价最高者赢得物品但只付第二高的价，此时``按真实估值出价''是每个人的占优策略——出高了可能以高于估值的价成交，出低了可能错失本该赚到的成交，而两种偏离都不会带来好处。VCG 机制把这个思路推广到多物品、多约束的分配问题：每个参与者支付的是自己的存在给别人造成的损失，于是说真话依然占优。这是一条定理，不是经验判断。

它的边界必须一并交代。VCG 一般不满足预算平衡，机制往往要贴钱或者留下一笔无法退还的余额。它的计算依赖于求解最优分配问题，而那个问题在组合拍卖里通常是 NP-难的，\hyperref[sec:09-Convex-Optimization/09-Complexity-and-Approximation/01]{``P、NP 与 NP-hard：给问题按计算难度分类''}里那条线在这里直接生效——机制在纸上成立，在规模上跑不动。它对合谋不稳健：几个参与者串通报价可以显著提高各自所得。收入也可能低得离谱，在某些配置下卖家几乎收不到钱，这是许多现实拍卖不用 VCG 的直接原因。

更硬的一条是 Myerson 与 Satterthwaite 的不可能定理：在买卖双方各自持有私有估值的双边交易里，没有任何机制能同时做到激励相容、双方自愿参与、预算平衡、以及事后有效率。四个要求里必须放弃一个。这与本部分其他几个单元的结论是同一类东西——不是还没找到好机制，是可以证明它不存在。

\subsection{哪些是定理，哪些是判断}

这一节的内容分两层，混起来会出问题。

Braess 例子里的算术是定理：给定那四条延迟函数与 \(4000\) 的流量，\(65\) 与 \(80\) 是逐条验证过的均衡值，不依赖任何建模假设之外的东西。仿射延迟下 \(\mathrm{PoA}\le4/3\) 是定理。第二价格拍卖与 VCG 的激励相容性是定理。Myerson--Satterthwaite 是定理。

``拆掉一条路可以改善交通''不是定理，是这个例子的一个特例。真实路网里是否存在 Braess 式的边、拆掉之后是否真的改善，要具体算，而且需求本身会随通行时间变化（路更快了就有更多人开车出门），这一层反馈不在上面的模型里。

``收费能改善拥堵''是有条件的分析结论。它成立的前提是需求对价格有弹性、存在可替代的出行方式、收费的执行开销不吞掉收益、以及分布效应可以接受——同一笔费用对不同收入的人不是同一件事，而模型里的``社会成本''把所有人的时间按同样的权重相加了。这个加总方式本身是一个价值判断，不是数学结论，\hyperref[sec:18-Synthesis/01-System-Lifecycle/01]{``把诉求写成目标函数与约束''}讲的正是这一步的分量。

``机制设计通常比劝导有效''是经验规律。它在很多场合被反复观察到，但没有定理保证，而且反例不少：规则一旦公开就会被针对，参与者会去优化被写下来的那个量而不是它本想代表的东西。\hyperref[sec:13-Machine-Learning/01-Learning-Problems-and-Evaluation/02]{``损失、风险、目标与指标不是同一个量''}与\hyperref[sec:18-Synthesis/01-System-Lifecycle/06]{``上线之后系统开始改变它所预测的世界''}讲的是同一件事的两个侧面。机制设计把``参与者会钻空子''从缺陷变成了前提，可它没有、也不可能保证写下来的规则就是我们真正想要的。

到这里，均衡这个概念本身的性质基本清楚了：\hyperref[sec:17-Constrained-Guarantees/06-Games-and-Equilibria/02]{``Nash 均衡存在但可能不唯一也不好''}给出存在性与两条坏消息，这一节把``不好''量化并给出改动的方向。剩下一个从头悬着的问题没有回答：均衡是一个静态条件，它只说``到了那里没人想走''，没说参与者怎样到达那里，甚至没说他们会不会到达。下一节把上一个单元的无遗憾算法接进来，看看两个学习者互相对打时，动力学真正收敛到的是什么。
